The entire arithmetic target can be described explicitly. Put
\(U=\Zh^\times\), let \(a,b\) be the quadratic characters of
\(\Q(\sqrt2)\) and \(\Q(i)\), and for each odd prime \(p\) put
\[
 p^*=(-1)^{(p-1)/2}p,\qquad
 c_p=\epsilon_{p^*},\qquad
 \varepsilon_p=\frac{1-(2/p)}2,\qquad
 \ell_p(d)=\frac{1-(d/p)}2.
\]
Here the last two quantities lie in \(\F_2\), and \(p\nmid d\)
in the last definition. Let \(\delta_2=b^2\). For odd \(p\),
let \(\delta_p\in H^2(U,\F_2)\) be the inflation of the central
cyclic extension
\[
 C_{2^{v_2(p-1)+1}}\longrightarrow C_{2^{v_2(p-1)}}
\]
from the maximal two-primary quotient of \(\Z_p^\times\).

\begin{theorem}\label{thm:main-full-quadratic-target}
With \(A=G/N=H/(H\cap N)\), transgression is an isomorphism
\[
 \operatorname{tr}:Y_2\xrightarrow{\ \sim\ }
 \ker\bigl(H^2(A,\F_2)\longrightarrow\operatorname{Br}(\Q)[2]\bigr).
\]
There is a direct-sum decomposition
\[
 Y_2=\F_2\eta_3
 \oplus\bigoplus_{p\text{ odd}}\F_2\psi_p
 \oplus\bigoplus_{p\text{ odd}}\F_2\kappa_{2,p}
 \oplus\bigoplus_{\substack{p<q\\p,q\text{ odd}}}\F_2\kappa_{p,q},
\]
where \(\eta_3=m_3/8\bmod2\), with
\(m_3(g)=[S^2T]f_g\) in the expansion \(x=1+S,y=1+T\),
is the cubic character and
the additional basis characters are uniquely specified by
\begin{align*}
 \operatorname{tr}(\kappa_{2,p})
 &=\operatorname{Inf}_U^A
       \bigl(ac_p+\varepsilon_p(\delta_2+\delta_p)\bigr),\\
 \operatorname{tr}(\kappa_{p,q})
 &=\operatorname{Inf}_U^A
       \bigl(c_pc_q+\ell_p(q^*)\delta_p+\ell_q(p^*)\delta_q\bigr).
\end{align*}
All products are cup products, and every character is a finite sum.
The character \(\kappa_{5,29}\) is realized by the central quadratic
layer of
\[
 \Q\bigl(\sqrt5,\sqrt{29},\sqrt{7+2\sqrt5}\bigr)
       /\Q(\sqrt5,\sqrt{29}),
\]
restricted to \(G_{K_2}\). It is nonzero and independent of
\(\eta_3\) and the \(\psi_p\).
\end{theorem}

Theorem~\ref{thm:fr-Y2-complete} proves the decomposition by
computing the affine cohomology and cancelling each cup product's
local Brauer invariants. The span of \(\eta_3\) and the
\(\psi_p\) is therefore a proper subspace of \(Y_2\), with quotient
of countably infinite dimension. The next two constructions lift all these basis directions to
\(X_2\), giving the complete restriction image.

For odd coefficients the translation coordinate disappears from the
cohomology, and the target has a uniform description. Fix an odd prime
\(p\), put \(\mathcal S_p=\{\ell:\ell\text{ prime},\ \ell\equiv1\pmod p\}\),
and choose cyclic coordinates in
\[
 U^{(p)}=\Z_p\times\prod_{\ell\in\mathcal S_p}C_{p^{v_p(\ell-1)}}.
\]
Write \(\gamma_0,\gamma_\ell\) for their reductions to \(\F_p\),
and \(\delta_\ell\) for the central cyclic extension class obtained
by increasing the exponent of the \(\ell\)-factor by one, with
kernel coordinate \(x\mapsto p^{v_p(\ell-1)}x\).

\begin{theorem}\label{thm:main-full-odd-target}
For every odd prime \(p\), inflation
\(H^r(U,\F_p)\to H^r(A,\F_p)\) is an isomorphism in every degree.
The arithmetic images of the \(\delta_\ell\) form a basis of
\(H^2(H,\F_p)\). Put
\(a_\ell=(\ell-1)/p^{v_p(\ell-1)}\bmod p\), and define
\begin{align*}
 D_{\ell m}&=a_\ell^{-1}\gamma_m(\ell)\delta_\ell
                  -a_m^{-1}\gamma_\ell(m)\delta_m\quad(\ell<m),\\
 D_{0\ell}&=-a_\ell^{-1}\gamma_0(\ell)\delta_\ell.
\end{align*}
Here Frobenius acts by the positive power map.
Transgression identifies \(Y_p\) with
\[
 \ker\bigl(H^2(A,\F_p)\longrightarrow H^2(H,\F_p)\bigr),
\]
and its inverse images of
\(\gamma_i\smile\gamma_j-D_{ij}\), for
\(i<j\) in \(\{0\}\cup\mathcal S_p\), form a basis of \(Y_p\).
Every character is a finite sum of these basis elements.
\end{theorem}

Theorem~\ref{thm:odd-full-galois-target} proves this by taking the
cyclotomic eigenspace of the global Brauer exact sequence over
\(\Q(\mu_p)\). Only primes congruent to one modulo \(p\) contribute,
and each cup product needs a correction at at most its two positive
indices. This determines the codomain of the odd-prime restriction
problem; its image and kernel for full \(G\) remain open.
